\documentclass[12pt]{article}
\usepackage[utf8]{inputenc}
\usepackage[T1]{fontenc}
\usepackage[margin=1in]{geometry}
\usepackage{setspace}
\usepackage{xcolor}
\usepackage[colorlinks=true, urlcolor=blue!60!black, linkcolor=blue!60!black]{hyperref}
\setlength{\parindent}{0em}
\setlength{\parskip}{0.6em}
\renewcommand{\baselinestretch}{1.0}
\begin{document}
\begin{center}
{\LARGE \textit{Regents of the University of California v. Bakke}, 438 U.S. 265 (1978)\par}
\vspace{0.5em}
{\large Case Brief --- Two-Page Edition\par}
\end{center}

\section{Memory Jogger}
UC Davis's 16-seat racial quota is struck down and Bakke admitted, but race may be one ``plus'' factor in admissions (\href{https://supreme.justia.com/cases/federal/us/438/265/}{Justia}).

\section{Detailed Case Facts}
UC Davis Medical School reserved 16 of 100 seats for a special program screened by a separate committee; its applicants, who identified as disadvantaged or minority, were exempt from the GPA cutoff and never compared with regular applicants. Allan Bakke, a white applicant with strong scores, was rejected in 1973 and 1974 while special admittees with substantially lower averages were admitted.

\section{Procedural History}
The California courts held the program unlawful, barred any use of race, and ordered Bakke admitted. The U.S. Supreme Court affirmed in part and reversed in part.

\section{Judicial Votes}
4--1--4 with six opinions. Powell joined Stevens, Burger, Stewart, and Rehnquist to strike the quota, and joined Brennan, White, Marshall, and Blackmun to permit some use of race. White, Marshall, and Blackmun also wrote separately.

\section{Holding}
Title VI tracks the Equal Protection Clause; a rigid quota is unlawful and Bakke must be admitted; but universities may consider race as one factor. Powell alone found diversity a compelling interest, a view \textit{Grutter} later adopted.

\section{Analysis of Opinions}
\textbf{Powell:} All racial classifications receive strict scrutiny because equal protection rights are personal. Remedying ``societal discrimination'' is too amorphous, but diversity, rooted in academic freedom, is compelling; Davis's quota was not necessary, while Harvard's ``plus'' approach was. \textbf{Brennan bloc:} benign remedial classifications warrant intermediate scrutiny, which the program met. \textbf{Stevens bloc:} Title VI is ``crystal clear'' in barring exclusion by race. \textbf{Blackmun:} ``In order to get beyond racism, we must first take account of race.''

\section{Examples: Future Applications}
\textbf{Same-Side:} an automatic 20-point bonus for minority applicants, as in \textit{Gratz}, is invalid. \textbf{Opposite-Side:} a single-pool holistic review treating race as a flexible ``plus'' is valid under \textit{Bakke}. \textbf{Fence-Sitter:} adding race as a tiebreaker to a race-neutral Top 10\% plan, as in \textit{Fisher}.

\section{Critique}
Antonin Scalia's \textit{The Disease as Cure} attacked racial preferences as perpetuating resentment, a view that prevailed in \textit{SFFA} (2023). Antisubordination scholars counter that Powell's symmetrical strict scrutiny made remediation of racial caste nearly impossible.

\section{Key Quotations}
\begin{quote}
``The guarantee of equal protection cannot mean one thing when applied to one individual and something else when applied to a person of another color. If both are not accorded the same protection, then it is not equal.''
\end{quote}
\noindent --- Justice Powell, announcing the judgment of the Court, 438 U.S. at 289--290.

\end{document}
